\documentclass[11pt,a4paper]{article}

% ---------- Packages ----------
\usepackage[margin=1in]{geometry}
\usepackage{amsmath, amssymb, amsthm, mathtools}
\usepackage{microtype}
\usepackage{algorithm}
\usepackage{algpseudocode}
\usepackage[colorlinks=true,
            linkcolor=blue!50!black,
            urlcolor=blue!50!black,
            citecolor=blue!50!black]{hyperref}

% ---------- Theorem environments ----------
\theoremstyle{plain}
\newtheorem{theorem}{Theorem}[section]
\newtheorem{proposition}[theorem]{Proposition}
\newtheorem{lemma}[theorem]{Lemma}
\newtheorem{corollary}[theorem]{Corollary}

\theoremstyle{definition}
\newtheorem{definition}[theorem]{Definition}

\theoremstyle{remark}
\newtheorem{remark}[theorem]{Remark}
\newtheorem{example}[theorem]{Example}

% ---------- Math shortcuts ----------
\newcommand{\R}{\mathbb{R}}
\newcommand{\N}{\mathbb{N}}
\newcommand{\E}{\mathbb{E}}
\newcommand{\Var}{\operatorname{Var}}
\newcommand{\F}{\mathcal{F}}
\newcommand{\Prob}{\mathbb{P}}
\newcommand{\Q}{\mathbb{Q}}
\newcommand{\dd}{\mathop{}\!\mathrm{d}}
\newcommand{\eqdef}{\overset{\mathrm{def}}{=}}

% ---------- Document metadata ----------
\title{Monte Carlo Pricing of a European Call:\\
       Milstein Discretisation under Geometric Brownian Motion}
\author{Quant Finance Portfolio --- Phase 2, Block 1, Algorithm 3}
\date{\today}

\begin{document}
\maketitle

\begin{abstract}
We specialise the Milstein scheme of the previous theoretical writeup
to geometric Brownian motion, derive the resulting recursion, and
apply it to Monte Carlo pricing of the European call. The principal
purpose of the algorithm is not improvement of the pricer itself ---
Milstein and Euler share weak order $1$ and produce statistically
indistinguishable European prices --- but rather \emph{empirical
demonstration of the strong/weak convergence asymmetry}. Running both
schemes on the same grid of step sizes and on the same Brownian
paths, we observe in a single experiment that Milstein achieves
strong order $1$ where Euler achieves only $\tfrac{1}{2}$, while the
two are tied at weak order $1$. This is the most direct possible
confirmation of the theoretical result, and it is the experimental
heart of Phase~2 Block~1.
\end{abstract}

\tableofcontents

\section{Introduction}
\label{sec:intro}

The previous writeup of Block~1.2.1 implemented Euler-Maruyama and
verified empirically that its strong order is $\tfrac{1}{2}$ and its
weak order is $1$. The asymmetry is a striking prediction of the
SDE convergence theory developed in Block~1.2.0: the same scheme,
applied to the same dynamics, has different convergence rates
depending on which kind of error one measures.

The natural follow-up question is whether higher-order schemes lift
both rates simultaneously, or only one. The answer, predicted by the
theory, is that the Milstein scheme lifts the strong order from
$\tfrac{1}{2}$ to $1$ but leaves the weak order unchanged at $1$. The
present writeup verifies this prediction empirically on the same
benchmark (European call pricing under GBM with all its inputs at
canonical values) used in Block~1.2.1, with the additional design
choice that the Euler and Milstein experiments share their step-size
grid so that the two slopes can be plotted on a single log-log axis.

The pricing application of the algorithm is, frankly, a sideshow.
Both schemes give weak order $1$ for European pricing under GBM, and
the closed-form sampler of Block~1.1 is unbiased and faster than
either; Milstein adds nothing to the pricing pipeline that Euler did
not already provide. The value of the algorithm is in the convergence
study, where the asymmetry is demonstrated in code rather than asserted
by a theorem. The practical relevance of Milstein in this project will
arise in Block~4 (pathwise sensitivity), where the strong order
matters because the simulated path must be a faithful pathwise
approximation of the true process for path-derivative estimators to
make sense.

\section{Milstein for geometric Brownian motion}
\label{sec:setting}

\subsection{The recursion}

Recall the GBM SDE under the risk-neutral measure $\Q$:
\begin{equation}
\dd S_t \;=\; r S_t \dd t + \sigma S_t \dd W_t,
\qquad S_0 > 0,
\end{equation}
so $a(s) = r s$ and $b(s) = \sigma s$, hence $b'(s) = \sigma$ and
$b(s) b'(s) = \sigma^2 s$. The Milstein recursion of the previous
writeup specialises to
\begin{equation}
\label{eq:milstein-gbm}
\hat S_{n+1} \;=\;
\hat S_n + r \hat S_n h + \sigma \hat S_n \Delta W_n
+ \tfrac{1}{2} \sigma^2 \hat S_n \big[(\Delta W_n)^2 - h\big],
\qquad \hat S_0 = S_0,
\end{equation}
with $\Delta W_n \sim \mathcal{N}(0, h)$ independent across $n$ and
$h = T / N$. Factoring $\hat S_n$ out:
\begin{equation}
\label{eq:milstein-gbm-mult}
\hat S_{n+1} \;=\; \hat S_n \cdot
\Big(\,1 + r h + \sigma \Delta W_n
+ \tfrac{1}{2} \sigma^2 \big[(\Delta W_n)^2 - h\big]\Big).
\end{equation}
The recursion remains multiplicative, but the multiplicative factor
is now \emph{quadratic} in $\Delta W_n$ rather than linear as in
Euler. This has two practical consequences worth noting.

\subsection{Improved positivity}

The Euler factor $1 + r h + \sigma \Delta W_n$ is unbounded below
along the negative direction: a sufficiently negative $\Delta W_n$
produces a negative factor and hence a negative $\hat S_{n+1}$. The
Milstein factor adds the term $\tfrac{1}{2} \sigma^2 ((\Delta W_n)^2
- h)$, which has \emph{no upper bound on the negative side}: the
$(\Delta W_n)^2$ component is always nonnegative, and only the $-
\sigma^2 h / 2$ piece pulls the factor down, by an amount bounded by
$\sigma^2 h / 2$. Concretely, the Milstein factor is bounded below
by
\begin{equation}
1 + r h + \sigma \Delta W_n - \tfrac{1}{2} \sigma^2 h,
\end{equation}
which is the worst case where $(\Delta W_n)^2 = 0$. For typical
finance parameters this lower bound is essentially the Euler factor,
so the positivity behaviour of Milstein is no worse than that of
Euler. Empirically, the Milstein-$\sigma^2/2$ correction more than
compensates the loss of the $\sigma \Delta W_n$ term in most
realisations: the factor $1 + r h + \sigma \Delta W_n + \tfrac{1}{2}
\sigma^2 ((\Delta W_n)^2 - h)$ is, on average, closer to $1$ than the
Euler factor.

We do not implement positivity defences in the algorithm; for the
parameter regime of our experiments, neither Euler nor Milstein
produces nonpositive prices in any realisation we have observed.

\subsection{The corrector term and its expectation}

The new term $\tfrac{1}{2} \sigma^2 ((\Delta W_n)^2 - h)$ is the
Milstein corrector. Two of its features deserve special note.

\paragraph{It has zero mean.} Since $(\Delta W_n)^2$ has mean
$\Var(\Delta W_n) = h$, the corrector has expectation zero:
$\E[\tfrac{1}{2} \sigma^2 ((\Delta W_n)^2 - h)] = 0$. This is the
mathematical reason that Milstein and Euler tie at weak order $1$:
the corrector affects each path differently, but its average
contribution to $\E[f(\hat S_T)]$ is, to leading order, zero.

\paragraph{It is pathwise nonzero.} Although the corrector vanishes
on average, on any individual path $(\Delta W_n)^2 - h$ takes a
specific value, typically of order $h$ in absolute terms. This is
the term that, in Euler, was discarded and contributed $O(h)$ to the
strong error per step. In Milstein, it is included, and the next
omitted term is of higher order: the strong error becomes $O(h)$
globally instead of $O(\sqrt{h})$.

The two features together explain the asymmetry. Strong convergence
sees the pathwise contribution of the corrector and benefits from
its inclusion. Weak convergence sees only the average contribution,
which is zero either way. This is the conceptual content of the
strong/weak distinction made concrete on a specific recursion.

\subsection{Convergence orders}

By Theorem~6.1 of the previous writeup applied to GBM,
\begin{equation}
\label{eq:milstein-orders}
\E[|S_T - \hat S_T^{\mathrm{Milstein}}|] \;=\; O(h),
\qquad
|\E[f(S_T)] - \E[f(\hat S_T^{\mathrm{Milstein}})]| \;=\; O(h)
\quad \text{for } f \text{ smooth enough}.
\end{equation}
Compare with Euler:
\begin{equation}
\E[|S_T - \hat S_T^{\mathrm{Euler}}|] \;=\; O(h^{1/2}),
\qquad
|\E[f(S_T)] - \E[f(\hat S_T^{\mathrm{Euler}})]| \;=\; O(h).
\end{equation}
Strong order: $\tfrac{1}{2}$ (Euler) vs $1$ (Milstein) --- a genuine
factor-$\sqrt{h}$ improvement at small $h$. Weak order: $1$ (Euler)
vs $1$ (Milstein) --- no improvement.

\section{The biased estimator and its comparison with Euler}
\label{sec:estimator}

\subsection{Definition}

Let $\hat S_T^{(i), \,\mathrm{Milstein}}$ for $i = 1, \ldots, M$ denote
independent realisations of the Milstein-discretised terminal price.
The Monte Carlo Milstein estimator of the European call price is
\begin{equation}
\hat C_{M, h}^{\mathrm{Milstein}} \;\eqdef\; \frac{1}{M} \sum_{i=1}^{M}
e^{-rT}\, (\hat S_T^{(i), \,\mathrm{Milstein}} - K)^+.
\end{equation}
The bias-variance decomposition of Block~1.2.1, equation~(8), applies
here verbatim: total mean square error decomposes as variance/$M$
plus bias-squared, with bias $O(h)$.

\subsection{Comparison with Euler at fixed budget}

The pricing performance of Milstein and Euler is, in this setting,
indistinguishable. Both have weak order $1$, so both produce a bias
of order $h$ for the same $h$. The constants in front of the $h$
asymptote may differ slightly between the two, but for canonical GBM
parameters the difference is too small to matter at the sample sizes
where the bias is detectable.

The cost per path of Milstein is marginally higher: one extra
multiplication and one squaring per Euler step (the corrector
$(\Delta W_n)^2 - h$). The overhead is on the order of $20$--$30\%$
in our implementation, well below the cost of generating the
$\Delta W_n$ themselves. So at fixed total budget, Milstein delivers
roughly the same accuracy as Euler with marginally fewer paths or
marginally larger $h$. Either way the difference is in the noise.

The genuine practical disadvantage of Milstein for European pricing,
visible in the validation suite of Section~\ref{sec:validation}, is
none whatsoever for European pricing as such: this whole comparison
is academic. The advantage materialises in pathwise applications.
We mention it here because confronting the irrelevance of Milstein
for vanilla European pricing is part of internalising the strong/weak
distinction.

\section{Empirical validation strategy}
\label{sec:validation}

The validation is the experimentally most informative part of the
algorithm. We design a single comparison study that exhibits both
the strong-order improvement and the weak-order tie of Milstein
relative to Euler.

\subsection{Test 1: side-by-side strong order}

The most distinctive part of the validation script is a comparison
of the two schemes' strong errors on the \emph{same} grid of step
counts and using the \emph{same} Brownian paths. The protocol:

\begin{enumerate}
\item Choose a logarithmic grid of step counts $N_k$ matching the
strong-order grid of Block~1.2.1, $N \in \{10, 30, 100, 300, 1000,
3000\}$.
\item For each $N_k$, sample $M$ matrices of Brownian increments
$\Delta W$ of shape $(M, N_k)$.
\item Run the Euler scheme on these increments to obtain $\hat S_T^{
\mathrm{Euler}, (i)}$ and the Milstein scheme on the \emph{same}
increments to obtain $\hat S_T^{\mathrm{Milstein}, (i)}$.
\item Compute the exact terminal $S_T^{(i)}$ from $W_T^{(i)} = \sum_n
\Delta W_n^{(i)}$ via the closed-form GBM solution.
\item Estimate
$\mathrm{strong}_{\mathrm{scheme}}(h_k) \approx M^{-1} \sum_i
|S_T^{(i)} - \hat S_T^{\mathrm{scheme}, (i)}|$
for both schemes.
\item Plot both error curves on a single log-log axis.
\end{enumerate}

The expected result is two parallel lines on log-log: Euler with
slope $\tfrac{1}{2}$, Milstein with slope $1$. Beyond confirming the
two convergence orders, the side-by-side plot makes the asymmetry
visually obvious: at the smallest $h$ tested, Milstein's strong error
should be roughly $\sqrt{h_{\mathrm{smallest}} / h_{\mathrm{largest}}}$
times smaller than Euler's, an order of magnitude or more.

\subsection{Test 2: weak order under common random numbers}

A direct adaptation of the Block~1.2.1 weak test, applied to
Milstein. The expected slope is $1$, the same as Euler. This test
serves as the empirical confirmation that Milstein \emph{does not}
improve the European pricer: same weak order, same asymptotic
quality.

The grid restriction discussed in the Block~1.2.1 numerical lesson
applies here: the bias must dominate the residual MC noise of the
CRN-corrected difference. With $M = 200{,}000$ and $N \in \{10, 20,
50, 100\}$, both schemes resolve the slope cleanly.

\subsection{Tests 3 and 4: coherence and validation}

Identical in structure to the corresponding tests of Block~1.2.1,
applied to the Milstein pricer. No conceptual novelty.

\section{Algorithmic specification}
\label{sec:algorithm}

\begin{algorithm}[H]
\caption{Milstein Monte Carlo for the European call (GBM)}
\label{alg:mc-milstein}
\begin{algorithmic}[1]
\Require Spot $S_0$, strike $K$, maturity $T$, rate $r$, volatility
$\sigma$, step count $N$, sample size $M$, confidence level
$1 - \beta$, random generator \texttt{rng}.
\Ensure Estimate $\hat C$, half-width $h_M$, sample variance $S_M^2$.
\State $h \gets T / N$,
$\quad D \gets e^{-rT}$
\For{$i = 1, \ldots, M$}
\State $S^{(i)} \gets S_0$
\For{$n = 0, 1, \ldots, N-1$}
\State Draw $\Delta W \sim \mathcal{N}(0, h)$ using \texttt{rng}
\State $S^{(i)} \gets S^{(i)} \cdot \big(1 + r h
+ \sigma \Delta W + \tfrac{1}{2} \sigma^2 ((\Delta W)^2 - h)\big)$
\EndFor
\State $Y_i \gets D \cdot \max(S^{(i)} - K,\, 0)$
\EndFor
\State $\hat C \gets M^{-1} \sum_i Y_i$
\State $S_M^2 \gets (M-1)^{-1} \sum_i (Y_i - \hat C)^2$
\State $h_M \gets z_{\beta/2}\, \sqrt{S_M^{\,2}/M}$
\State \Return $(\hat C,\, h_M,\, S_M^{\,2})$
\end{algorithmic}
\end{algorithm}

The implementation vectorises the inner loop across paths (advance
all $M$ states simultaneously) using \texttt{numpy.cumprod} on the
matrix of factors, identical in structure to the Block~1.2.1 Euler
implementation. The only changes to the code are the addition of the
quadratic corrector term in the factor and the corresponding
docstrings; the Block~1.2.2 implementation is barely 30 lines longer
than Block~1.2.1.

\section{Bibliographic notes}
\label{sec:biblio}

The Milstein scheme is from Milstein (1974) and is treated in detail
in Kloeden \& Platen~\cite{KloedenPlaten}, Section~10.3. The empirical
side-by-side comparison study designed here is a direct adaptation
of the comparison study in Higham~\cite{Higham2001}, Section~5. For
the practical relevance of Milstein in finance applications,
Glasserman~\cite[Section~6.3]{Glasserman} provides the
context-specific discussion: when does the strong order of a scheme
matter for the financial applications at hand.

\begin{thebibliography}{99}

\bibitem{Glasserman}
P.\ Glasserman,
\emph{Monte Carlo Methods in Financial Engineering},
Springer, 2004.

\bibitem{Higham2001}
D.\ J.\ Higham,
\emph{An algorithmic introduction to numerical simulation of
stochastic differential equations},
SIAM Review, 43(3):525--546, 2001.

\bibitem{KloedenPlaten}
P.\ E.\ Kloeden and E.\ Platen,
\emph{Numerical Solution of Stochastic Differential Equations},
Springer, 1992.

\end{thebibliography}

\end{document}
